
\section{Tutorial}
\label{sec:manual-tutorial}

Welcome to CubedOS!

\subsection{Introduction}
\label{sec:manual-tutorial-introduction}

This document is for new CubedOS developers who want to learn how to build CubedOS applications.
Here we will walk you through the process of creating a simple application and then refer you to
the \filename{samples} folder in the CubedOS GitHub repository for additional examples.

CubedOS applications consist of one or more \newterm{domains} where each domain consists of one
or more \newterm{modules}. Simple (and even many not-so-simple) applications are built as a
single domain. Multi-domain applications are outside the scope of this tutorial, but the
Networking sample shows an example of such an application.

Each module in a domain communicates with other modules by passing messages. Sometimes modules
act as servers; they receive request messages, process those requests, and return reply messages
to the requester. Sometimes modules act as clients; they send request messages to servers,
receive the replies, and do other processing. It is normal for modules to play both roles at
different times; they receive request messages and send requests of their own as they process
their incoming requests.

CubedOS encourages a client/server application architecture, or, more generally, a distributed
application architecture with messages being sent between modules freely. However, the
Request/Reply terminology is conventionally used throughout. In some cases ``requests'' are sent
even when no reply is expected, and in other cases ``replies'' are sent unsolicited.

Normally modules are written in SPARK/Ada as a package hierarchy. The top-level, parent package
has the same name as the module. Child packages are used to implement the various subsystems
within a module. Some child packages have conventional names, as we will describe later. This
precise organization is not strictly required, but it is recommended and certain aspects of the
system assume this organization. It is possible, in principle, to write CubedOS modules in other
languages such as full Ada or even C. We do not discuss how to do that in this tutorial.

Communication in CubedOS is point-to-point, asynchronous, and unreliable. There is no built-in
mechanism for broadcasting or multicasting messages; doing so requires writing code that
explicitly sends multiple point-to-point messages as needed. There is, however, a
publish/subscribe core module that offers a ``channel'' abstraction. Modules can subscribe to
particular channels and receive messages that are published to those channels by other modules.
This decouples senders and receivers so they are not aware of each other and provides a form of
multicasting.

Messages are asynchronous in the sense that the sender is allowed to continue before the message
has been received. The send operation only queues the message for later delivery. After a
successful send, the sending module can reuse any data structures it populated for the outgoing
message. Every module has at least one task, called the \newterm{module main task} that receives
and processes messages. Modules may have other internal tasks for other purposes. Multiple
modules execute in parallel. It is possible, even common, for senders to send multiple requests,
possibly to different modules, before getting any replies, and replies may arrive in a different
order from which the requests were sent. It is up to you to coordinate this activity in your
modules, although CubedOS provides some assistance, as we will discuss.

Messages are unreliable in the sense that final delivery is not guaranteed. Senders can elect to
receive a status code to inform them if their messages are dropped by the message passing
infrastructure, but even this status code does not indicate if the message was actually received
by the destination module (for example, if the destination module is stuck or crashed or
unreachable over a network). If a server module wishes to indicate an error in a request, or any
other kind of failure, it must send an explicit reply message saying so.

Every CubedOS domain contains an array of ``mailboxes'' that hold pending messages. There is a
one-to-one correspondence between the modules of a domain and the mailboxes. In other words,
each module has exactly one mailbox and each mailbox services exactly one module. The main task
of a module is normally blocked trying to read a message from that module's mailbox. When a
message is available, the main task fetches it, processes it, and then (likely) blocks again
trying to read the next message. Since the mailboxes have, in general, space for holding
multiple messages, it is possible that a message can arrive while the module was processing an
earlier message. In that case, the new message is retrieved at once. Modules never process more
than one message at a time. Outstanding messages are queued in the module's mailbox. The mailbox
array thus defines the domain and forms the communication infrastructure for the domain.

In normal operation, all module tasks are blocked, and the CPU is idle. When a hardware
interrupt is generated, an interrupt service routine inside one a driver module for the
associated hardware device activates and sends a message, as appropriate, to another,
higher-level module to handle the interrupt. That module may send other messages, etc., creating
a cascade of activity in the system. Once the interrupt has been fully handled and that activity
dies down, all modules return to their quiescent state of being blocked waiting for a message.
If multiple interrupts happen at about the same time, this presents no problem as the modules
are able to queue messages and execute in parallel as needed. However, the modules do need to be
coded with this scenario in mind. CubedOS is thus a \newterm{reactive} programming system since
it reacts to external input, but does not typically initiate any activity.

CubedOS comes with a collection of \newterm{core modules} that provide services commonly needed
by many applications. The publish/subscribe server mentioned earlier is an example of a core
module. There is also a timeserver module that can be requested to send messages to other
modules at specified times or periodically. "Reacting" to such messages gives modules a way to
perform activities without other specific hardware stimulation (i.e., routine monitoring,
telemetry reporting, and so forth).

CubedOS also comes with a collection of \newterm{library packages} that contain general-purpose
reusable code of various kinds. Unlike a module, a library package does \emph{not} contain a
task\footnote{In a previous version of CubedOS, library packages were called ``passive modules''
to distinguish them from active modules with tasks. That terminology has been dropped.}.
Subprograms inside library packages are called in the usual way and return results in the usual
way, and not via messages\footnote{It is permitted for a library package to send a message as a
side effect.}.

Library packages can also be shared by modules. One consequence of this is that library packages
can have no modifiable global data\footnote{Unless wrapped in a protected object.} since they
might be called by multiple tasks at once. All task interactions are via the message passing
mechanism. Modules should not attempt to communicate via global data inside library packages.

The rest of this document walks you through the steps for creating a simple, single-domain
CubedOS application. You can use these steps as a framework for creating your own applications.
Although we build the application here incrementally for educational purposes, you can review
the final form of the application in the \filename{moonshot} folder of the \filename{samples}
folder in the repository.

\subsection{Moonshot Application}
\label{sec:manual-tutorial-moonshot}

The application we will build demonstrates many, but not all, aspects of CubedOS development. It
is intended to be easy to understand, implement, and modify. It is not intended to be an
exhaustive demonstration of CubedOS and its libraries.

We will create an application, called \newterm{Moonshot}, that roughly simulates a CubeSat
mission orbiting the Moon. It is not our intention for this application to be an accurate
simulation of a real mission. Many details will be left out or simplified. However, imagining
that the application might be controlling a real spacecraft will hopefully make it more fun to
think about!

Although in a real mission, the application would obviously need to control real hardware, in
this tutorial we will simulate the hardware. The modules that might ordinarily serve as device
drivers will instead implement ``mock'' devices with simulated behavior. This allows the
application to run on any development system you might have available.

Moonshot will make use of several CubedOS core modules. Specifically:
\begin{enumerate}
  \item The Log Server will be used to log messages of various kinds for later transmission to
  the Earth. The Log Server is widely used. Almost all CubedOS applications will require it.

  \item The Time Server will be used to trigger periodic actions in other modules. In Moonshot,
  without any real hardware to generate interrupts, the Time Server will be the initiator of all
  activity once system initialization is complete.

  \item The File Server will be used to store image files taken by the mock camera. The Moonshot
  mission will entail taking picture of the lunar surface, so we need a place to store those
  images. In a real mission it is common for there to be some sort of non-volatile storage
  precisely for this purpose.

  \item The Publish/Subscribe Server will be used to allow modules to subscribe to various
  channels. The point of the Publish/Subscribe Server is to decouple senders and receivers of
  messages so that reusable modules can be more readily created. We describe this in more detail
  later.
\end{enumerate}

Moonshot will also have several application-specific modules for managing the (mock) hardware
devices. Specifically:
\begin{enumerate}
  \item The Camera module will simulate a camera that can take pictures of the lunar surface.

  \item The Radio module will simulate a radio that can send messages to and receive messages
  from the Earth.

  % \item A Thruster module will simulate maneuvering thrusters that can adjust the orbit and
  % orientation of the CubeSat.

  \item A Controller module that will server as the main control logic for the mission. It will
  coordinate the action of the other modules to fulfill the mission objectives.
\end{enumerate}

To get started, you will need space for your application's files. You will also need to clone
the CubedOS repository from GitHub. At the time of this writing CubedOS is distributed in source
form via its development repository. There is no release distribution yet.

The folder for your application need not be included in the CubedOS repository folder. In fact,
we recommend that you do \emph{not} include it there. Instead, create a folder somewhere else on
your system for your application.

It is often useful to use an existing application as a starting point for new work. The Echo
sample in the CubedOS repository is a good reference. You may wish to look try running this sample
and look back at it while developing the Moonshot tutorial. First, create a \filename{build} and
\filename{src} folder in your application folder. Also, create a folder \filename{src/mxdr} as a
placeholder for the MXDR files you will create later. The precise folder structure you use isn't
mandated by CubedOS, but this is what we recommend for CubedOS applications.

Next, create a GNAT project file in the root of your application folder named
\texttt{moonshot.gpr} that contains the following:
\begin{verbatim}
project Moonshot is
   for Main use ("main.adb");
   for Object_Dir use "build";
   for Source_Dirs
      use ("src", "src/mxdr", "../../src/modules", "../../src/library");
   for Languages use ("Ada");

   package Compiler is
      for Default_Switches ("Ada")
        use ("-gnat2022",     -- Use Ada 2022
             "-gnatW8",       -- Enable UTF-8 source files
             "-fstack-check", -- Enable stack checking
             "-gnatwa",       -- Enable all warnings
             "-gnata",        -- Enable assertions
             "-g");           -- Enable debugging
   end Compiler;
end Moonshot;
\end{verbatim}

The source paths should include two source folders from the CubedOS repository. Edit the paths
above as appropriate for your system. The default switches are fine for this tutorial
application. They enable all warnings, runtime assertion checking, and debugging support. In a
real application you may want different switches or multiple build scenarios.

Notice that CubedOS applications are intended to use Ada 2022 as their base language. Although
this is not strictly necessary, Ada 2022 does provide some niceties that are helpful. Finally,
the |-gnatW8| switch enables support for Unicode UTF-8 encoded source files, allowing certain
Unicode characters, such as Greek letters, in identifier names. That is also not strictly
necessary, but it can be helpful for clarity in some cases.

Every Ada program requires a library-level, parameter-less procedure to serve as the program's
entry point. CubedOS applications don't actually use this procedure, but it is required for the
Ada language. Use the name \texttt{main.adb} as indicated above. We will describe what must be
in this file later.

Filling in the rest of the application requires the following steps:
\begin{enumerate}
\item Instantiating the CubedOS generic message manager to create a domain (i.e., a mailbox
  array) for your application.
\item Writing a name resolver package to assign addresses to the modules you will use, including
  the core modules.
\item Creating skeletons for the modules you will write by copying the templates from the
  CubedOS repository and editing them.
\item Writing the Ada-required main procedure to ensure the modules you need will be compiled
  into your executable.
\item Verifying that your completed skeleton application builds and runs.
\item Writing MXDR interface descriptions for the messages that will be sent between modules,
  and using the Mercury tool to generate the Ada code for those message encoding and decoding
  subprograms.
\item Writing a publish/subscribe channel definition package to define the channels you will use
  in your application.
\item Implementing the full logic of your application by filling in the rest of your
  application's details.
\end{enumerate}

\subsection{Instantiating the Message Manager}
\label{sec:manual-tutorial-instantiating}

Every domain in CubedOS requires a message manager to handle the message passing between
modules. The message manager is a generic package that is instantiated with several arguments
that define the number of modules in the domain and the dimensions of the mailboxes used by
those modules.

At the time of this writing, the generic instance \emph{must} be called |Message_Manager|. This
name is hard-coded in the core modules. In a multi-domain application, you would have multiple
instances of the generic message manager, one for each domain, but since every instance must
have the same name, they must reside in different executable units. This is a limitation of the
current CubedOS design that might be lifted in the future, if it proves to be overly
restrictive.

Here is a suitable instantiation of the generic message manager for the Moonshot application:

\begin{verbatim}
  with CubedOS.Generic_Message_Manager;
  pragma Elaborate_All(CubedOS.Generic_Message_Manager);

  package Message_Manager is
    new CubedOS.Generic_Message_Manager
      (Domain_Number =>  1,
       Module_Count  =>  8,
       Mailbox_Size  =>  8,
       Maximum_Message_Size => 1024);
\end{verbatim}

The meaning of the generic arguments is as follows:

\begin{description}
  \item[Domain\_Number] The domain number is a unique identifier for the domain. It is a
  positive integer. For a single domain application, such as Moonshot, the domain number should
  always be 1.

  \item[Module\_Count] The module count is the number of modules in the domain. For the Moonshot
  application, we have four application-specific modules (Camera, Radio, Thruster, and
  Controller) and four core modules (Log Server, Time Server, File Server, and Publish/Subscribe
  Server). That makes a total of eight modules.

  \item[Mailbox\_Size] The mailbox size is the number of messages that can be queued in a
  module's mailbox (also called the ``depth'' of the mailbox). The mailbox size should be large
  enough to hold all the messages that might be sent to a module before it can process them.

  It can be difficult to decide the most appropriate value for this parameter. If a module can
  always process messages quickly, or if messages arrive infrequently, a small mailbox size is
  appropriate. However, if a module can be slow to process messages or if messages can sometimes
  arrive quickly, a larger mailbox size is needed. A mailbox size that is too small can lead to
  messages being dropped. However, a large mailbox size wastes memory.

  Be aware that all modules in the same domain will use the same mailbox size. Choose a size
  large enough to satisfy the needs of the module requiring the greatest size of any in the
  domain. A mailbox size of 8 is a reasonable default.

  \item[Maximum\_Message\_Size] The maximum message size is the maximum number of 8-bit bytes
  (also called ``Octets'') that can be sent in a single message. The maximum message size should
  be large enough to hold the largest message that will ever be sent in the domain.
\end{description}

The total amount of memory (not counting certain overheads) required for the mailboxes is the
product of |Module_Count|, |Mailbox_Size|, and |Maximum_Message_Size|. In the Moonshot
application, the total memory required for the mailboxes is $8 \times 8 \times 1024 = 64$ KiB.
Depending on the platform on which the application will run, this may be a significant amount of
memory. Be sure to consider the memory constraints of your platform when choosing these
parameters.

\subsection{Writing the Name Resolver}
\label{sec:manual-tutorial-name-resolver}

In CubedOS, every module has an address represented by a pair of integers (Domain ID, Module
ID). In the Moonshot application, the domain ID is 1, and the module IDs are in the range from 1
to 8 inclusive.

You must assign these module IDs to the modules yourself, in any way you see fit, but you must
ensure that each module has a unique ID. The core modules do \emph{not} have predefined ID
numbers. You must assign them as well.

The ID assignments are made in a package called |Name_Resolver|. As with the message Manager
package, the name resolve package must also have a specific name because references to it are
hard coded in the core modules. Here is a suitable |Name_Resolver| package for the Moonshot
application:

\begin{verbatim}
with Message_Manager; use Message_Manager;

package Name_Resolver is

   -- Core Modules
   Log_Server              : constant Message_Address := (0, 1);
   Time_Server             : constant Message_Address := (0, 2);
   File_Server             : constant Message_Address := (0, 3);
   Publish_Subscribe_Server: constant Message_Address := (0, 4);

   -- Application-Specific Modules
   Camera                  : constant Message_Address := (0, 5);
   Radio                   : constant Message_Address := (0, 6);
  %  Thruster                : constant Message_Address := (0, 7);
   Controller              : constant Message_Address := (0, 8);

end Name_Resolver;
\end{verbatim}

The |Message_Address| type is defined in the |Message_Manager| package. It is a record type with
two components for the domain ID and module ID.

Note that, although all the modules in our application are in domain 1, the value 0 is used for
the domain ID in the |Name_Resolver|. This value is treated as ``this'' domain by the message
passing system. If a module sends a message to domain 0, it is sending that message to a module
in the same domain as the sender. The explicit value of 1 would also work in this case, but
using 0 is a bit more general. The use of domain ID 0 is more important in a multi-domain
application when you want to send a message to a module in the same domain without knowing your
own domain ID.

It is important that you assign unique module IDs to each module in your application. If two
modules are given the same ID, they will both try to fetch messages from the same mailbox.
Although, it is permitted under the Jorvik profile used by CubedOS for multiple tasks to queue
on the same mailbox entry, it is unspecified which messages will be picked up by which tasks.
This will result in the two conflicting modules receiving many messages intended for the other,
producing many ``invalid message'' errors at runtime.

It is also important that you don't accidentally assign a module ID outside the range of allowed
ID values implied by the |Module_Count| parameter of the |Message_Manager| instantiation.
However, in this case you will get compile-time warnings, so it is less likely to be a problem.

Finally, there is no requirement that the module IDs be contiguous or start at 1. For example,
you could assign the core modules IDs 1 through 4 and the application-specific modules IDs 10
through 13, leaving the range 5 through 9 reserved for future expansion without renumbering the
existing modules. You may notice that the module ID of 7 is missing in the above example. This
is because in a future version of this tutorial, we may want other modules (such as a thruster).
Leaving a gap allows us to fill the module in, without having to renumber the existing ones.
However, this does cause the mailboxes to use more memory than is necessary, since you
are allocating mailboxes for modules you aren't using. Also, renumbering modules later in
development is not necessarily a problem since the compiler will manage the module ID values for
you when it rebuilds your application.

At the time of this writing, module ID 2 is reserved for the inter-domain router used in
multi-domain applications. This router is internal to CubedOS and not something application
developers interact with directly. In a multi-domain application, module ID 2 cannot be used for
either a core module or an application module. However, in a single domain application such as
Moonshot, module ID 2 is available for use.

\subsubsection{A Word About Addresses}
\label{sec:manual-tutorial-addresses}

As you have seen modules have two-dimensional addresses in CubedOS, consisting of domain ID,
module ID pairs. Each module also has a collection of messages that it can send and receive, and
those messages have IDs as well. The message IDs are local to the module. The module IDs are
local to the domain. The domain IDs are globally unique.

The full address of a message is a triple consisting of the domain ID, module ID, and message
ID. For example: $(1, 3, 2)$ is the address of message 2 in module 3 of domain 1. This is a
different message from $(1, 3, 3)$ in the same module, and different from $(1, 4, 2)$ in a
different module in the same domain, and of course different from $(2, 3, 2)$ in a different
domain altogether. In most cases, you do not need to work with message IDs directly. The CubedOS
message passing system takes care of that for you. However, you should be aware of this
addressing scheme in case you need to debug a problem or understand the CubedOS internals.

\subsection{Creating Module Skeletons}
\label{sec:manual-tutorial-skeletons}

Now we are ready to start building the modules for the Moonshot application. To start,
we will focus on the controller module. The controller module is the main control logic for the
mission, and thus serves as the ``main'' program for the application.

CubedOS modules are organized as a package hierarchy. The top-level package has the same name as
the module, and is mostly used as a place to hang various child packages that implement the module's
logic. Here is a complete top-level package for the controller module:

\begin{verbatim}
pragma SPARK_Mode(On);

package Controller is
  -- Empty
end Controller;
\end{verbatim}

This package only serves to form the root of the module's package hierarchy, and thus does not
need to contain any functionality (or have a body). It might be reasonable to include some
module-wide declarations here, such as types needed throughout the module, but in this case we
have none.

Recall that in Ada, child packages have visibility into their parents automatically (meaning
without using |with| statements). The bulk of the module can be in child packages, or even in
grandchild packages, and yet can see the module-global declarations in the top-level package
without making those declarations global program-wide. This is a nice example of Ada's scope and
visibility-management features in action.

CubedOS modules normally have at least two child packages: |API| and |Messages|.

The |API| package contains subprograms for creating messages from well-typed values, breaking
messages back into well-typed values, and testing messages for validity. The API of a module is
all about which messages that module receives (called ``Requests'') and sends (called
``Replies'').

Characterizing messages as requests and replies follows from the view of message passing as a
kind of client/server interaction where the client makes a request, and the server replies to
it. However, in CubedOS, modules can act as both clients and servers, depending on what they are
doing. Also, requests can be sent without expecting a reply, and replies can be sent without ever
receiving a request (i.e., unsolicited replies).

The |Messages| package contains a task that fetches a message from the module's mailbox,
processes that message, and then fetches the next message. This task is called the module's main
task.

The CubedOS repository includes a folder \filename{templates} that contains template files for
both the |API| and |Messages| packages. You can copy the templates you need into your module's
\filename{src} folder and edit them to suit your needs. They are heavily commented to help guide
you. The expectation is that you will ultimately remove most, if not all, of the provided
comments and replacement them with application-specific comments as appropriate.

In most cases, however, as we describe in Section~\ref{sec:manual-tutorial-mercury}, you would
use the \newterm{Mercury} tool to generate a module's |API| package for you and not rely on the
templates for that package.

The |Messages| package template shows the recommended structure of that package. The
implementation of the main task (at the bottom) is very simple. It immediately passes each
message it receives to procedure |Process|. That procedure, in turn, examines the message and
calls the appropriate subprogram to handle that message.

Handling a message entails using an API decoding procedure to first break the message into its
well-typed components, and then acting on the message. Normally, this involves sending other
messages to other modules, and ultimately, in most cases, sending a reply message back to the
requester.

Copy the template |Messages| package templates to your \filename{src} folder, renaming the files
as appropriate for the Controller module. Edit the templates to change names, as necessary. For
now, don't try to implement any real logic. Let's get the skeleton to compile and run first.

\subsection{Writing the Main Procedure}
\label{sec:manual-tutorial-main-procedure}

All Ada programs require a library-level main procedure. In a CubedOS application, this
procedure does nothing, but it is where we will inform the compiler of the modules we want to
include in the executable, so it knows what code to locate and compile. Here is a suitable main
procedure for the Moonshot application so far:

\begin{verbatim}
-- Bring in the necessary modules.
with CubedOS.Log_Server.Messages;
with CubedOS.Time_Server.Messages;
with CubedOS.File_Server.Messages;
with CubedOS.Publish_Subscribe_Server.Messages;
with Controller.Messages;
with Camera.Messages; -- we'll make this later, so it doesn't have to be included yet

pragma Unreferenced(CubedOS.Log_Server.Messages);
pragma Unreferenced(CubedOS.Time_Server.Messages);
pragma Unreferenced(CubedOS.File_Server.Messages);
pragma Unreferenced(CubedOS.Publish_Subscribe_Server.Messages);
pragma Unreferenced(Controller.Messages);
pragma Unreferenced(Camera.Messages);

procedure Main is
begin
   null;
end Main;
\end{verbatim}

Although |Main| returns immediately, the process won't terminate until all the module tasks have
terminated. However, the module tasks are all infinite loops (as required by the Jorvik profile)
so the program never ends. This is common for embedded systems. The control program runs for as
long as the machine it is controlling runs.

The |with| statements bring the module tasks into the executable. The Ada compiler will form the
transitive closure of package dependencies, locating all the necessary program components, based
on these |with| statements and on |with| statements in the included packages, etc. Since no
module directly interacts with the message loop of any other module, it is essential that you
|with| the |Messages| packages here. Otherwise, a module's message loop will not be linked into
your executable and that module's mailbox will never be serviced.

The |pragma Unreferenced| directives suppress warnings about unused |with| statements, but have
no other effect.

Notice that we are including all the core modules we will use, but only the Controller
skeleton module we made previously. We will add the other application-specific modules as we
create them, editing |Main| to include them when we are ready.

At this point, the program should compile and run. When executed, all modules will block
immediately trying to fetch a message from their mailboxes. Since no messages are ever sent,
nothing more happens. We're off to a good start!

\subsection{Using Mercury}
\label{sec:manual-tutorial-mercury}

If you spent much time studying the template API packages, you probably noticed that they are
intricate. Each encoding and decoding subprogram involves repeated calls to the XDR library to
encode (or decode) individual message parameters into (or from) an array of raw octets. This is
tedious and error-prone to write manually.

To help with this, CubedOS includes a tool called \newterm{Mercury} that generates the encoding
and decoding subprograms for you. Mercury reads a simple, human-readable description of the
messages sent and received by a module, and it outputs a SPARK-provable API package for that
module.

Mercury is a work in progress, and it may not be able to handle all your API needs
automatically. However, even when the API package it produces is incomplete or insufficient, you
can still use Mercury as a starting point. It is much easier to make adjustments to the
generated code than to write the entire API package from scratch.

Like all message-passing systems, CubedOS is inherently distributed. In principle, messages can
be passed from one module to another in the same process, in different processes on the same
computer, or even on different computers connected by a network. From the point of view of one
module, all other modules are ``remote'' code. Multi-domain CubedOS applications capitalize on
this inherently distributed nature to allow applications to be spread across a network of
spacecraft in a natural way.

Many distributed application environments exist: COBRA, ONC RPC, SOAP, Ice, and gRPC, to name
only a few. \pet{Citations needed} Essentially all of these environments use some
\newterm{interface definition language} (IDL) to describe the interface of remote code. That IDL
is then compiled into client-side code (sometimes called ``stubs'') that \newterm{marshal} the
parameters of a remote procedure call into a message that can be sent over a network, and
server-side code (sometimes called ``skeletons'') that \newterm{unmarshal} the message back into
parameters that can be passed to the actual remote code. These stubs and skeletons, together
with the message-passing infrastructure connecting client and server form the
\newterm{middleware} layer of the application.

In CubedOS terms, the middleware layer is the message manager together with the API packages
generated by Mercury. Marshaling is done by the encoding functions generated by Mercury, and
unmarshaling is done by the decoding functions generated by Mercury. Because the API packages
are part of the middleware layer provided by the system, it is unreasonable and inappropriate
for application developers to be required to write them manually. Thus, Mercury is not just a
convenience; it is an essential part of CubedOS's architecture.

CubedOS currently uses XDR (eXternal Data Representation) \cite{rfc-4506} as its encoding
format, but it could just as easily use some other approach such as ASN.1, Protocol Buffers, or
even JSON. \pet{Citations needed} As a CubedOS application developer, switching to a different
encoding format should be as simple as setting an option on Mercury's command line. You should
not be required to manually rewrite the marshaling and unmarshaling code in your application
(i.e., the API packages).

\subsubsection{MXDR}

The interface definition language used by Mercury is an extension of XDR \cite{rfc-4506} called
MXDR (Modified XDR). MXDR includes features for specifying constraints on basic types, and for
specifying messages and their direction of information flow. We will illustrate MXDR using the
Moonshot application.

Consider, for example, the Camera module. For simplicity, we will assume that it only knows
about two messages: a request to take an image, and a reply containing the file name of where
it stored the image in the file system. It would be reasonable for the Camera module to send
the image data directly, but in this case it replies with what is, in effect, a pointer to
where that data can be found. Here is the MXDR description of these two messages:

\begin{verbatim}
message struct -> Take_Image_Request {
    void;
};

message struct <- Take_Image_Reply {
    string File_Name<128>;
};
\end{verbatim}

The core of an MXDR description are the message definitions. Messages are conceptualized as
structures with named fields for the message parameters. The direction of the arrow indicates
the direction of information flow. An arrow pointing to the right is for messages that go from a
client to the module being described. An arrow pointing to the left is for messages that go from
the module being described to the client. Thus, an MXDR description is always from the point of
view of the module being described as playing the role of a server. When a module acts as a
client, it will use the API packages of the other modules that service its requests.

In this case the Camera module provides an operation ``Take Image'' and so accepts a message
|Take_Image_Request| and replies with a message |Take_Image_Reply|. The names of the messages
are arbitrary, but the use of ``Request'' and ``Reply'' is conventional. The names do not
necessarily need to be related, i.e., the reply message could be named in some other way if
desired.

In this simple example, the request messages have no parameters. In a more realistic example, it
is likely that additional parameters would be needed to specify exactly how the camera should
capture the image. Those additional parameters would be fields in the message structure. The
special field value |void| is used here to explicitly indicate that there are no other
parameters. It is required to be used in such cases.

The reply message contains a single parameter, a string of up to 128 characters. The string type
maps to the usual Ada string type, but in this case, there will be an additional decoded
value to indicate how much of the string is actually used.

The MXDR description is saved in a file with the extension \filename{.mxdr}, in this case
\filename{Camera.mxdr} (with the indicated casing), and should be included in the
\filename{src/mxdr} folder of your application. The Ada code generated by Mercury will also be
saved to \filename{src/mxdr}. This separates the interfaces of your modules from their
implementations, which is a good practice. It is similar to the way some C and C++ projects use
a separate folder for header files.

Mercury uses the name of the MXDR file to construct the name of the generated Ada package. In
this example, Mercury will create the package |Camera.API| in the files
\filename{camera-api.ads} and \filename{camera-api.adb}. The package will contain the following
six subprograms:

\begin{verbatim}
  function  Take_Image_Request_Encode(...) return Message_Record;
  function  Is_Take_Image_Request(...) return Boolean;

  function  Take_Image_Reply_Encode(...) return Message_Record;
  function  Is_Take_Image_Reply(...) return Boolean;
  procedure Take_Image_Reply_Decode(...);
\end{verbatim}

You may notice that there is no |Take_Image_Request_Decode|. This is because, as you may
remember, the request messages have no parameters. Thus, there is nothing to decode. Mercury
can detect this and will make no decode function. There is however, an encode function. This
is because, even though there is no XDR encoding of parameters, an empty message still must be
constructed with the correct sender address and request ID.

For clarity, the details of the subprogram parameters are left out. However, all CubedOS
messages include a |Request_ID| parameter that the sender can use to mark messages with An
identifier meaningful to that sender. This is useful for correlating requests and replies.

For example, the sender might command the camera to take several images, perhaps using different
parameters. The sender can mark each request with a unique |Request_ID|. When the Camera module
replies, it will copy the sender's |Request_ID| into the reply. The sender can then use that
information to match the reply with a particular request. If there is no need for a sender to
correlate requests and replies, the |Request_ID| can be set to zero and ignored.

Keep in mind that the sender might send multiple requests before receiving any replies, allowing
it to tend to other matters while the camera acquires the images and writes them to storage.
This highly asynchronous behavior is typical of CubedOS applications.

To understand how to use a module, you only need to read the MXDR for that module, resorting to
the generated API package only when you need to know the exact subprogram signatures. For
example, if the Controller module wants to access the image file created by the Camera module,
it will need to send messages to the file server module. Here is a part of the file server's
MXDR:

\begin{verbatim}
enum Mode_Type { Read, Write };

// The special value 0 is used for invalid files.
typedef unsigned int File_Handle_Type
    range 0 .. 64;
typedef File_Handle_Type Valid_File_Handle_Type
    range 1 .. File_Handle_Type'Last;

// Types to count the number of octets read/written.
const Maximum_Read_Size  = 256;
typedef unsigned int Read_Result_Size_Type
    range 0 .. Maximum_Read_Size;
typedef unsigned int Read_Size_Type
    range 1 .. Read_Result_Size_Type'Last;

// Open the named file in the given mode.
message struct -> Open_Request {
    Mode_Type Mode;
    string    Name<>;
};

// Returns the invalid handle if the open operation failed.
message struct <- Open_Reply {
    File_Handle_Type Handle;
};

// Request that Amount octets be read from the indicated file.
message struct -> Read_Request {
    Valid_File_Handle_Type Handle;
    Read_Size_Type         Amount;
};

// Returns data from the file.
message struct <- Read_Reply {
  Valid_File_Handle_Type Handle;
  Read_Result_Size_TYpe  Amount;
  opaque                 File_Data<>;
};
\end{verbatim}

The protocol speaks for itself. First a module must send an |Open_Request| message to the file
server, specifying the desired file name. The server responds with an |Open_Reply| message
containing a ``handle'' to the file (or zero if the open failed). The module then sends repeated
|Read_Request| messages to the file server specifying the previously opened handle. The server
responds with a |Read_Reply| messages containing the requested data.

The full interface to the file server also includes messages for writing files, closing files,
and some other operations. The file server implementation provided with CubedOS uses the file
system on your development system. In a real mission, the implementation (i.e., the |Messages|
package) could be replaced with a version that works with whatever storage is available on the
spacecraft without changing the interface or the modules that rely on that interface.

Notice how various types are defined before the message structures can be declared. Mercury will
generate corresponding Ada type definitions into the API package accordingly. As an extension to
XDR, Mercury allows you to specify constrained ranges on the types, following the usual Ada
style. This is essential. Mercury attempts to generate SPARK-provable code, and SPARK's proofs
are greatly facilitated by the use of constrained types. This feature sets Mercury apart from
most other IDL compilers.

\subsubsection{Invoking Mercury}

Mercury is included in the CubedOS repository in the \filename{mercury} folder. It must be built
before it can be used. Fortunately, this is easy to do. However, since Mercury is written in
Scala, it may require that you install some additional tools. Specifically:

\begin{enumerate}
\item You must install a Java runtime. Mercury requires Java 21 or later. Download Java from
  \url{https://www.oracle.com/java/technologies/downloads/}.

\item You must also install SBT. Download SBT from \url{https://www.scala-sbt.org/}. Any recent
  version of SBT will be fine. It fetches and caches the precise version required by a project
  when you first build that project. SBT is a build tool for Scala, similar to Maven or Gradle
  for Java projects.
\end{enumerate}

Significantly, you do not need to explicitly install Scala. SBT will fetch the correct version
of Scala for you, along with any other required libraries and plug-ins, when you build Mercury
for the first time.

Next, open a command prompt or terminal session, and change to the \filename{mercury} folder of
the repository. Execute the command \command{sbt assembly}. This will build Mercury and create a
standalone executable JAR file containing the tool. Although not required, you might also want
to use the command \command{sbt test} to run the test suite for Mercury.

There is a shell script named \filename{mercury.sh} in the \filename{bin} folder off the root of
the repository. This script is a convenience wrapper around the Java command to run Mercury. A
symbolic link to this wrapper script can be created wherever you find convenient.

Mercury requires template files that it uses during code generation. These templates are located
in the \filename{mercury/templates} folder of the repository. The \filename{mercury.sh} script
specifies this location to Mercury as a command line option. Unless you want to use different
templates (which you almost never will), you do not need to worry about this.

Once Mercury is ready, apply the shell script wrapper to your MXDR file to generate the API
packages as follows:

\begin{verbatim}
  $ ./mercury.sh Camera.mxdr
\end{verbatim}

Note that this only produces \filename{camera-api.ads} and \filename{camera-api.adb}. It does
not produce the top level package itself (\filename{camera.ads} in this case). When configuring
the build of your application, you must arrange for Mercury to be run on the MXDR files, if
necessary, before the Ada compiler is invoked.

During normal development, you will most likely want to integrate the execution of Mercury
into your regular build environment. In a realistic application, with many modules (and thus
many MXDR files), you will want the build system to recompile only the MXDR files that have
changed, if any, since the last build. This can be accomplished using the \command{make}
utility. Here is a \filename{Makefile} that can be added to the \filename{mxdr} folder of
your Moonshot application.

\begin{verbatim}
# Adjust to reflect Mercury's location.
MERCURY=../../../../bin/mercury.sh

# List of generated specification files.
all:	camera-api.ads

.PHONY: clean

clean:
        rm -f *.ads *.adb

camera-api.ads:	Camera.mxdr
        $(MERCURY) $^
\end{verbatim}

This file only describes how to build the Ada specifications; the body files are ignored. Since
Mercury generates both files together, the \command{make} utility only needs to focus on the
specifications, letting the bodies be generated as a kind of side effect.

To add a new MXDR file, you only need to add the name of the specification to the list of
prerequisites to the ``all'' target, and include a suitable rule for the new target file at the
bottom.

For full effect, you now only need to integrate \command{make} in your overall build system.
\pet{Is there a way to get GNAT Studio to run \command{make} automatically as part of the
build?}

Mercury is intended to produce files that can be provided free of runtime errors with \SPARK.
Thus, you can (and should) use the \SPARK\ tools to verify the generated code along with your
handwritten code. If a bug in Mercury causes it to generate buggy code, you will want to give
SPARK an opportunity to find that error!

Also, the files produced by Mercury are intended to be human-readable and human-editable.
Ideally you would not need to edit them, but it might be necessary in some cases. Those cases
could reasonably be described as exhibiting a deficiency in Mercury, but the tool is still
useful even when it cannot completely automate the generation of the API packages in all
situations.

There is much more that can be said about MXDR and using Mercury. We will describe additional
features as we encounter them in the Moonshot application. A complete tutorial for Mercury is in
Section~\ref{sec:manual-tutorial-mercury}.

\subsection{Moonshot MXDR}
\label{sec:manual-tutorial-moonshot-mxdr}

Now we are ready to write the MXDR for the application-specific modules in the Moonshot application.
In a real application we would likely have detailed requirements to help guide us in this task.
However, in this tutorial we will use highly simplified interfaces since we are only trying to
illustrate the application development process.

\subsubsection{Camera Module}

The MXDR for the Camera module presented in Section~\ref{sec:manual-tutorial-mercury} is good
enough for our purposes. We repeat it here for convenience:

\begin{verbatim}
  // FILE: Camera.mxdr

  message struct -> Take_Image_Request {
      void;
  };

  message struct <- Take_Image_Reply {
      string File_Name<128>;
  };
\end{verbatim}

\subsubsection{Radio Module}

\todo{Finish Me!}

% \subsubsection{Thruster Module}

% \todo{Finish Me!}

\subsubsection{Controller Module}

The MXDR for the Controller module is unusual in that there is no MXDR file required! The
Controller module plays the role of the ``main'' program of the application. It coordinates the
activity of the rest of the system. As a result, it only acts like a client to the other modules
and provides no services that the other modules will use. It accepts no request messages, and
has no corresponding replies. It therefor has no MXDR interface definition and no API package.

\subsection{Writing a publish/subscribe channel}
\label{sec:manual-tutorial-pub-sub}

...

\subsection{Writing the Application}
\label{sec:manual-tutorial-writing}

Now that we have the skeletons of the modules and the MXDR descriptions of the messages, we can
begin to write the application logic. Let's build this application from the bottom up, starting
with the Camera module.

\subsubsection{Camera Module}

We have only one message to implement, |Handle_Take_Image_Request|. In this tutorial we will always
return one of two hard-coded images. The file \filename{Image-Copernicus.jpg} is a picture of
the Copernicus crater \cite{Copernicus}, and the file \filename{Image-Korolev.jpg} is a picture
of the Korolev crater \cite{Korolev} (on the far side of the Moon).

In a real application this module would need to control the camera hardware, capture an image,
and store it in a file. For demonstration, however, we will just alternate between the two
images when a request is received, sending first the Copernicus image and then the Korolev
image.

Once we decide what image to return, we simply send a |Take_Image_Reply| message to the sender
using the |Route_Message| procedure provided by the CubedOS message manager. Normally, we would
need to decode the |Take_Image_Request| message to extract any desired parameters from it.
However, in this application the message is empty so no decoding is required. However, we will
still need to encode the reply using a suitable encoder from the API package generated by
Mercury.

\subsubsection{Radio Module}

\todo{Finish Me!}

% \subsubsection{Thruster Module}

% \todo{Finish Me!}

\subsubsection{Controller Module}

The Controller plays the role of the application's main program. It performs some initialization
duties and then coordinates the behavior of the rest of the system.

Since CubedOS applications are reactive\footnote{They react to external events, but otherwise
don't initiate any activity.}, we will describe the behavior of the Controller module, not in
procedural terms, but in terms of how it responds to events (i.e., messages).

In a typical application, interrupts drive the action of the system. When a hardware device
creates an interrupt, the module that serves as that device's driver will service the interrupt
and, in general, send messages to other modules to take additional action. Those modules will,
in turn, send messages to still other modules, resulting in a flurry of activity across the
system. Once the interrupt, and all its consequences have been fully processed, system activity
dies down again, and all modules are left blocked on their mailboxes waiting for the next
message.

In this tutorial application, there is no physical hardware to create interrupts. Instead, we
will use the timeserver to generate periodic ``tick'' messages to kick-start Controller
activity. This is a normal use-case for the timeserver. In a real application, the tick messages
can be used to trigger periodic housekeeping activities, backups, or (slow) polling of hardware
components.

During initialization, before entering the main message fetching loop, we do:

\begin{itemize}
\item Configure a \newterm{periodic series} with the time server, requesting tick messages every
minute. We will need to specify a \newterm{series ID} for this series. Any value will suffice
(it is up to us), but we should remember it. In general, an application might have multiple time
series active at once, and we will want to distinguish one from the next using their unique
series ID values.
\end{itemize}

In this simple application, there are no other initialization activities.

Table~\ref{tbl:controller-message-actions} shows the Controller's responses to various messages.

% This table structure came from ChatGPT. It seems to work well!
\begin{table}[ht]
    \centering
    \caption{Controller Message Actions}
    \label{tbl:controller-message-actions}
    \renewcommand{\arraystretch}{1.5} % Adjusts row height
    \begin{tabularx}{\textwidth}{|>{\hsize=0.35\hsize}X|>{\hsize=0.65\hsize}X|}
        \hline
        \textbf{Message} & \textbf{Action} \\
        \hline \hline
        Tick message from the time server. &

        The time server is trying to wake us up, so we can do something. As the Camera to take a
        picture. \\
        \hline
        Camera replies with image file name. &

        Ordinarily the Camera might alternatively reply with an error, but in this simple
        application that will never happen. Send a message to the file server to open the image
        file. \\
        \hline
        File server replies to open request. &

        If the open request failed, log the error. Otherwise, prepare a space packet with a data
        size of 512 octets and APID of 42 (an arbitrary value selected by us). Ask the file
        server to read the first 256 octets of file data. \\
        \hline
        File server replies with data. &

        Pack the data received into the currently active space packet. If the packet is full,
        send it to the Radio module for transmission to Earth and prepare a new space packet. If
        256 octets received, ask the file server for the next 256 octets, otherwise ask the file
        server to close the file. \\
        \hline
        File server replies to close. &

        If the currently active space packet is partially filled, send it to the Radio module
        for transmission. \\
        \hline
    \end{tabularx}
\end{table}

There is an implicit assumption in Table~\ref{tbl:controller-message-actions} that the
Controller will be able to completely handle one image file before the time server says it's
time to ask the Camera for another. In this simple tutorial application, where the Camera can
acquire an image ``instantly,'' that assumption is reasonable. In a more realistic application,
it might be necessary to plan for the possibility of there being multiple active images at once.
In that case, the Controller could store pending image file names in a data structure that is
global to the module. When one file is closed, the Controller could request that the next be
opened, thus serializing what would otherwise be interleaved operations.

CubedOS's message passing paradigm would also allow interleaved file reads by using different
request IDs for the different files. When the module receives some file data, it could consult
the request ID on the received message to determine to which data stream that data belongs. Of
course, the Radio module would have to also be prepared to accept interleaved file transfers.
These are some of the design questions that arise when building CubedOS applications.

In this tutorial application, the Controller repackages file data into \newterm{space packets}
as defined by the Consultative Committee on Space Data Systems (CCSDS) \pet{Citation needed}. It
is the space packets that are sent to the radio, encapsulated in CubedOS messages. The CCSDS is
clear that space packets can be used for intermodule communication, as well as for
spacecraft-to-spacecraft and spacecraft-to-Earth communication. The CubedOS library contains a
package, |CubedOS.Lib.Space_Packets|, that facilitates the construction and manipulation of space
packets.

The header of a space packet contains an \newterm{application ID} (APID), that applications can
use to identify the semantic significance of the packet's contents. In our case, we arbitrarily
use an APID of 42 for file data. This choice would need to be known to ground software so that
it can distinguish file data from other kinds of telemetry.

Note also, that space packets contain a flag that can be used to distinguish the first, last,
and continuation packets of a data stream that spans multiple packets. We will use that flag to
ensure that ground software understands that multiple packets are used to hold the data from a
single file.

It would be reasonable for the Controller to not dirty itself with low-level activities
like reading files. Instead, it could pass the name of the file to the Radio and let that
other module deal with the details of getting the file data. However, a reusable Radio module
would mostly likely only want to work with space packets, leading to the approach here.

Actually, in a real mission, it is most likely that an entirely separate module implementing the
CCDS File Deliver Protocol (CFDP) would take care of reading the file and building CFDP protocol
units to be encapsulated in space packets. The Controller would just pass the name of the image
file to the CFDP module for transfer.
